\section{Computational evidence}\label{sec:experiments}

This section reports two archived computational studies. Neither was rerun
for this paper, and their results are not pooled: they used different
solvers, budgets, cohorts, and questions. The first, a precursor study,
tested iterated OBBT as an external presolve. Its findings shaped the second,
a prospective comparison of the numerical policy of
Section~\ref{sec:algorithms-policy} with native SCIP. All numbers come from
the archived records, which the supplementary archive contains. Quantities
computed afterwards from the per-run records, rather than by the archived
analysis, are identified as such; they are descriptive and were not part of the
frozen protocol.

\subsection{A precursor study of external iterated OBBT}\label{sec:experiments-precursor}

The precursor study asked whether iterating OBBT before a solve pays off.
Its candidates were the 377 MINLPLib instances~\cite{bussieck2003-minlpliba-collection-of-test-models,vigerske2026-minlplib-a-library-of-mixed} with quadratic nonlinearities
that are not flagged convex, have 5--2000 variables and no semicontinuous or SOS
variables, and have a known optimal value. It excluded 38 of them because
feasibility-based bound tightening left some variable of a nonlinear term
without a finite bound of magnitude at most $10^7$, leaving 339 instances from
85 families. The
solvers were Gurobi~13~\cite{gurobi2026-reference-manual} and
SCIP~10~\cite{hojny2025-the-scip-optimization-suite-10}, with one thread each,
two random seeds, and a
600-second limit on total time. A pipeline first ran the solver at the root node
for at most 30 seconds to obtain an incumbent $\hat f$, set the cutoff
$U=\hat f+10^{-6}\max\{1,|\hat f|\}$, ran OBBT on a separate McCormick LP with
five-point square tangents and secants, and restarted the solver on the
tightened box with the incumbent and cutoff. OBBT used sequential coordinate
updates with feasibility filtering. Five stopping rules were compared: one
round, five rounds, a rule that stopped after a round without change or at a
limit of 50 rounds or 400 seconds, and two adaptive rules. The adaptive rules
computed, over the variables that moved in a round, the ratio $\rho$ of their
total normalized width after the round to that before it, and stopped once
$\rho$ exceeded $0.5$ or $0.8$, after a round without change, or after 120
seconds of OBBT; later rounds were restricted to moved variables and their
bilinear partners. A control arm repeated the root run and the restart without OBBT,
which separates the effect of the box from the effect of restarting with an
incumbent. Appendix~\ref{sec:precursor-details} gives the tables.

No stopping rule solved more instance--seed pairs than the default solver:
OBBT pipelines solved 637--638 of 678 pairs with Gurobi against 640 by default,
and 561--571 with SCIP against 575. Every pipeline had a larger shifted
geometric mean of total time than the default, by factors of 1.31--1.73 for
Gurobi and 1.21--1.32 for SCIP. These ratios, like the node and final-solve
ratios below, are the archived summaries; Appendix~\ref{sec:precursor-details}
explains why two solved counts are corrected here but not in these summaries.
On the instances where OBBT ran,
the tightened boxes reduced node counts against the control---to $0.65$
(Gurobi) and $0.62$ (SCIP) of the control with the no-change-or-limit rule---but
the time of the final solve was $0.96$--$1.02$ of the control's, so the savings
in the final solve did not offset the OBBT time. In the group of pairs where
that rule's box closed at least half of the gap of the McCormick LP, the node
count fell to about one fifth of the control, but the mean control time of this
group was only 1.2 seconds (SCIP) and 4.9 seconds (Gurobi). More rounds
strengthened the relaxation
with diminishing returns. With the known optimal value as cutoff, OBBT closed
on average $0.236$ of the McCormick LP gap after one round, $0.322$ after five,
and $0.346$ with the no-change-or-limit rule; with the root incumbents that
rule closed $0.250$ (Gurobi) and $0.287$ (SCIP). The median closure was zero for
every rule. Each of the 4,573 tightened boxes that belonged to an instance with
a reference solution contained that solution, up to a numerical tolerance of
$10^{-6}(1+|x_j|)$ per coordinate.

Two observations from this study shaped the integrated policy. First, an
external presolve pays for a second root and a restart, which by themselves
made the control slower than the default (shifted geometric mean ratios $1.08$
for Gurobi and $1.16$ for SCIP). An integrated propagator avoids both. Second,
much OBBT work was spent where it could not help: on 109 instances whose
nonlinear variables are all integer, the first round changed a bound on only 25,
yet these instances consumed more than nine tenths of all first-round OBBT time.
This suggested a short pilot before a full round, reconsideration after an
incumbent improves, and a test cohort that the default solver does not solve
quickly. The prospective comparison implemented these suggestions.

\subsection{Design of the prospective comparison}\label{sec:experiments-design}

The comparison asks whether adding the propagator of
Section~\ref{sec:algorithms-policy} to SCIP improves solving within a short
budget, and whether its adaptive variant reduces the propagator's cost
relative to its fixed-order variant. The protocol, the cohort, the policy
parameters, and the analysis were frozen before any comparative outcome was
observed.

\emph{Arms.} The \emph{native} arm runs SCIP with default settings. The
\emph{fixed} and \emph{adaptive} arms add the propagator with the parameters of
Table~\ref{tab:policy-parameters}; they share the relaxation, the bound
validation, the budgets, and the triggers, and differ in ranking, screening,
and the pilot. All arms use identical SCIP settings and the same model
construction: quadratic rows become SCIP nonlinear constraints, a quadratic
objective is moved into a constraint on an auxiliary variable, and maximization
is converted to minimization. In SCIP~10.0.2 the LP-based OBBT propagator runs at
the root node by default, and the separate nonlinear OBBT propagator is
disabled. Generalized variable bound propagation is enabled by
default~\cite{scip1002-obbt,scip1002-nlobbt,scip1002-genvbounds}; these defaults
were left unchanged in every arm. The comparison
therefore measures additional work on top of native OBBT, not the effect of
OBBT itself.

\emph{Cohort.} Table~\ref{tab:cohort} lists the twenty models. The twelve
public models were selected by a rule fixed in advance from the precursor
study's in-scope instances: 5--300 variables, 1--200 nonlinear variables, at
most 500 distinct quadratic terms, and at least five seconds in every
historical native SCIP run. Of the 41 eligible names, ordered by the SHA-256
digest of a fixed prefix followed by the name, the first twelve were taken, at
most one per historical family. No known optimal value or outcome of the
current arms entered the rule. The cohort is a selected set of harder
instances, not a random sample. Its family labels still separate closely
related applications: four of the twelve are circle- or point-packing models.
Three public models contain binary variables, and four have variables without
finite bounds in their original formulation. The eight synthetic models, two
sizes each of four structures generated with seed 20261003, are stress tests:
point packing, which maximizes the smallest squared distance among $k$ points
in the unit square; coupled square budgets over $[-1,1]^n$, with concave
diagonal and cyclic bilinear objective terms, two quadratic budget rows, and a
two-sided sum row; bilinear cycles, with products $x_ix_{i+1}$ of $[0,1]$
variables around a cycle and a two-sided sum row; and sparse indefinite
quadratic programs over $[0,1]^n$ with random two-sided linear rows. Four tiny
synthetic models of other sizes and another seed were used before freezing to
check the interface and callback coverage; they are excluded from every result.
The normalized public inputs were compared with the original OSiL files at five
test points per model, and every returned public incumbent was also evaluated
directly from the original files.

\begin{table}[htbp]
\centering
\footnotesize
\setlength{\tabcolsep}{4pt}
\begin{tabular}{@{}lrrrr@{\quad}lrrr@{}}
\toprule
Public model & $n$ & $m$ & nonlin. & hist.\ (s) & Synthetic model & $n$ & $m$ & nonlin.\\
\midrule
\texttt{bayes2\_50} & 86 & 77 & 65 & 600 & packing, $k=5$ & 11 & 10 & 10\\
\texttt{ex5\_2\_5} & 32 & 19 & 27 & 600 & packing, $k=7$ & 15 & 21 & 14\\
\texttt{gabriel01} & 215 & 467 & 100 & 371 & coupled squares, 12 & 12 & 3 & 12\\
\texttt{genpooling\_meyer04} & 118 & 141 & 28 & 600 & coupled squares, 24 & 24 & 3 & 24\\
\texttt{graphpart\_3pm-0344-0344} & 144 & 48 & 144 & 12 & bilinear cycle, 12 & 24 & 13 & 12\\
\texttt{kall\_circles\_c6b} & 18 & 54 & 16 & 73 & bilinear cycle, 20 & 40 & 21 & 20\\
\texttt{kall\_circlespolygons\_c1p13} & 43 & 48 & 17 & 7 & indefinite QP, 20 & 20 & 5 & 20\\
\texttt{kall\_circlesrectangles\_c1r12} & 49 & 52 & 19 & 15 & indefinite QP, 36 & 36 & 9 & 36\\
\texttt{pointpack08} & 17 & 36 & 16 & 25 & & & & \\
\texttt{powerflow0009r} & 60 & 103 & 57 & 20 & & & & \\
\texttt{qp3} & 100 & 52 & 100 & 600 & & & & \\
\texttt{wastewater15m2} & 133 & 151 & 24 & 63 & & & & \\
\bottomrule
\end{tabular}
\caption{The twenty models of the prospective comparison: variables $n$,
rows $m$, and variables occurring in quadratic terms. ``Hist.'' is the
shorter of the two native SCIP times in the precursor study, rounded; 600
means that both runs reached the 600-second limit.}
\label{tab:cohort}
\end{table}

\emph{Runs and resources.} Every model was run with SCIP random seed shifts $0$
and $1$ in all three arms: 120 runs, 40 per arm. The two seeds are repeated runs
of one model, not independent problem families. Each run had a ten-second limit
on the solver call, which includes model construction and all propagator work;
construction time is subtracted from SCIP's own limit. A driver started each run
in a fresh process, recorded the elapsed time of the whole process---interpreter
start, model reading, solving, validation, and output---and enforced a
30-second hard limit. SCIP, its LP solver, and the numerical libraries used one
thread. At most two runs executed at a time, in an order shuffled with seed
20261003, on a 36-logical-CPU machine that was also running unrelated jobs; the
load averages at the start were 17.76, 15.65, and 14.11. The software was SCIP
10.0.2 through PySCIPOpt 6.2.1~\cite{maher2016-pyscipopt-mathematical-programming-in-python},
with Python 3.13.11, NumPy 2.5.3, and SciPy 1.18.1, whose HiGHS dual simplex
solved the auxiliary LPs~\cite{huangfu2018-parallelizing-the-dual-revised-simplex}.

\emph{Outcomes.} Each returned incumbent was checked against the original
bounds, integrality, quadratic rows, and objective of the stored model. A row
residual is scaled by the larger of one plus the sum of absolute term values at
the point and the absolute right-hand side, and a bound residual by one plus the
absolute coordinate; the incumbent passes if the largest scaled residual and the
integrality residual are at most $10^{-6}$. The objective is evaluated exactly
in rational arithmetic before conversion. This is a numerical check, not an
exact certificate. A run counts as solved if SCIP reports optimality, the
incumbent passes the check, the process did not fail, the dual bound does not
exceed the incumbent value by more than $10^{-6}$ relative, and the reported and
independently evaluated objectives agree to $10^{-6}$ relative. The primary
time summary is the mean PAR-2 score over all 40 runs of an arm: the process time
of a solved run and $20$ seconds, twice the limit, otherwise. The shifted
geometric mean uses the uncapped process times of all runs with a shift of one
second. With $p$ the validated incumbent value and $d$ the final dual bound,
both in minimization form, the normalized gap is
$\max\{0,p-d\}/\max\{1,|p|,|d|\}$; the gap score is this value capped at one, or
one if there is no valid incumbent or no finite consistent dual bound. Paired
comparisons with native use the same model and seed. A PAR-2 win or loss
requires a value below $95\%$ or above $105\%$ of native's, and a gap win or loss
requires a difference in normalized gap larger than $10^{-4}$, among pairs in
which both gaps are defined.

\subsection{Outcomes}\label{sec:experiments-outcomes}

All 120 runs completed. There were no process failures, invalid incumbents,
objective mismatches, or inconsistent primal and dual bounds. All 117 returned
incumbents passed the check, with largest absolute residual
$9.97\times10^{-7}$ and largest scaled residual $7.21\times10^{-7}$. Each arm had
one run without an incumbent, \texttt{genpooling\_meyer04} with seed $0$.
Table~\ref{tab:campaign} gives the main summaries.

\begin{table}[htbp]
\centering
\small
\begin{tabular}{@{}llrrrr@{}}
\toprule
Cohort & Arm & Solved & Mean PAR-2 (s) & Shifted mean (s) & Mean gap score\\
\midrule
All & Native & 14/40 & 13.855 & 5.957 & 0.20351\\
All & Fixed & 14/40 & 13.988 & 6.285 & 0.18438\\
All & Adaptive & 14/40 & 13.997 & 6.265 & 0.18806\\
\addlinespace
Public & Native & 2/24 & 18.964 & 10.248 & 0.33408\\
Public & Fixed & 2/24 & 18.999 & 10.297 & 0.30209\\
Public & Adaptive & 2/24 & 19.026 & 10.324 & 0.30798\\
\addlinespace
Synthetic & Native & 12/16 & 6.192 & 2.384 & 0.00767\\
Synthetic & Fixed & 12/16 & 6.471 & 2.772 & 0.00782\\
Synthetic & Adaptive & 12/16 & 6.453 & 2.733 & 0.00818\\
\bottomrule
\end{tabular}
\caption{Outcomes of the prospective comparison. Times are process times in
seconds. The shifted mean is the geometric mean of uncapped process times with
a shift of one second. Lower gap scores are better.}
\label{tab:campaign}
\end{table}

\emph{Solved runs.} Every arm solved the same 14 runs: both seeds of
\texttt{kall\_circlespolygons\_c1p13} and 12 of the 16 synthetic runs. Neither
added policy gained or lost a solve. Because the solved sets coincide and every
unsolved run scores 20 seconds, the differences in mean PAR-2 come entirely from
the process times of these 14 runs. In total they took 34.21 seconds for native,
39.52 for fixed, and 39.88 for adaptive, increases of 5.32 and 5.67 seconds.
The propagator's recorded time in the same runs was 5.23 seconds (fixed) and
5.72 seconds (adaptive), comparable to these increases. In paired terms, mean
PAR-2 rose by 0.133 seconds (fixed) and 0.142 seconds (adaptive);
adaptive had two PAR-2 wins and eleven losses, fixed one win and ten losses, all
among the solved runs. Both adaptive wins are the two seeds of the bilinear
cycle with twelve factors. There both added arms finished at the root node,
whereas native SCIP needed seven nodes; the adaptive process times were 0.1
seconds below native's, while the fixed arm matched native's time in one seed.
The callback records of these runs show the incumbent trigger at work: in both
arms and both seeds, the first root callback, made before any incumbent existed,
accepted no reduction; the callback triggered by the first incumbent accepted
six, and the following callback, triggered by the resulting domain change, six
more. The driver waited for each worker with Python's standard polling loop,
whose interval grows to 50 milliseconds, and the recorded times of short runs
cluster in steps of about 0.05 seconds. The time outside the solver
call---interpreter start, reading, validation, output, and the wait---was
0.41--0.59 seconds per run.

\emph{Final gaps.} Both added policies lowered the mean public gap score, from
0.334 to 0.302 (fixed) and 0.308 (adaptive). The two seeds of
\texttt{bayes2\_50} account for this. There both added arms ended with an
incumbent of value 0.520 instead of 6.825 while the dual bound stayed near zero,
which changed the gap score from 1 to 0.520. The propagator constructs no
incumbents, so this is an indirect effect of a changed search. Without
\texttt{bayes2\_50}, the mean public gap score was 0.274 for native, 0.282 for
fixed, and 0.289 for adaptive; these values were computed afterwards from the
per-run records. Among the 39 pairs in which both gaps are defined, adaptive had four
gap wins and twelve losses, and fixed had five wins and twelve losses; 23 of
these pairs are public, with four (adaptive) and five (fixed) wins and ten losses
each. In eleven of the twelve losses of each arm, the final dual bound was weaker
than native's. Gaps at a wall-clock limit are sensitive to how much search fits
into the limit: on \texttt{qp3}, where the propagator could never run, the fixed
arm still ended one seed with a different dual bound from native.

\subsection{Where the added work went}\label{sec:experiments-work}

Table~\ref{tab:work} decomposes the propagator's work, using the per-callback
records of the archived runs, which store each callback's node number and
depth. A callback is attributed to the root node, to a nonroot node with a
positive number, or to a nonroot node with number zero. The release-tagged
SCIP~10.0.2 source initializes new nodes with number zero, assigns positive
numbers to search-tree children, and leaves temporary probing nodes at
zero~\cite{scip1002tree}. Consistently, in runs that finished at the root,
every nonroot callback had number zero. This supports interpreting them as
temporary probing nodes, although the logs did not record node types. We report
them as nonroot, number-zero callbacks.
Because the trigger state is keyed by node number,
the policy treated all nonroot nodes with number zero in a run as one node and allowed at most
three acting callbacks for them. The time in the row of triggers without a call is the
remainder of the propagator time: trigger evaluations that started no acting callback,
which are timed but produce no record.

\begin{table}[htbp]
\centering
\footnotesize
\setlength{\tabcolsep}{4pt}
\begin{tabular}{@{}lrrrrrrrr@{}}
\toprule
& \multicolumn{4}{c}{Fixed} & \multicolumn{4}{c}{Adaptive}\\
\cmidrule(lr){2-5}\cmidrule(l){6-9}
& Calls & LPs & Accepted & Time (s) & Calls & LPs & Accepted & Time (s)\\
\midrule
Root node & 100 & 1150 & 62 & 8.224 & 98 & 336 & 62 & 3.554\\
Tree nodes & 166 & 662 & 50 & 5.485 & 505 & 1203 & 244 & 13.287\\
Nonroot, number 0 & 90 & 360 & 0 & 3.684 & 90 & 180 & 2 & 2.593\\
Unsupported box (\texttt{qp3}) & 48 & 0 & 0 & 0.068 & 48 & 0 & 0 & 0.066\\
Triggers without a call & & & & 2.069 & & & & 3.623\\
\midrule
Total & 404 & 2172 & 112 & 19.530 & 741 & 1719 & 308 & 23.123\\
\addlinespace
\quad of which LP solves & & & & 12.818 & & & & 10.393\\
\bottomrule
\end{tabular}
\caption{Work of the added propagator over all 40 runs of each arm.
``Accepted'' counts bound changes that SCIP applied. The time of LP solves
includes the dual-bound validation of each solve. The totals are those of the
archived analysis; the split was computed afterwards from the per-callback
records.}
\label{tab:work}
\end{table}

The adaptive variant solved $1{,}719$ directional LPs against $2{,}172$, a
reduction of $20.9\%$, but its propagator time rose from 19.530 to 23.123 seconds,
an increase of $18.4\%$. The decomposition explains the difference. The time
spent in LP solves and dual-bound validation fell by 2.43 seconds, while the
remaining propagator time rose from 6.71 to 12.73 seconds. This remaining time,
including trigger visits that started no acting callback, amounted to about 17
milliseconds per recorded callback in both variants, and an LP with its
validation took about 6 milliseconds. The adaptive variant made
337 more callbacks; 574 of its 693 supported callbacks stopped after the
two-LP pilot. Its 64-LP budget was therefore rarely binding: 20 of its 38 runs
on models other than \texttt{qp3} ended at the cap of 24 callbacks, and only 5 of
its 40 runs used all 64 LPs. The fixed variant used all 64 LPs in 27 runs and
never reached the callback cap outside \texttt{qp3}; 349 of its 356 supported
callbacks ended at an LP limit. Much of the reduction in LPs is thus a consequence
of the callback cap binding before the LP budget, not of evidence that the
adaptive variant chose its directions better. The same budgets also kept the
recorded work small: the largest propagator total of a run was 1.00501 seconds
against an allowance of one second, 7 fixed and 11 adaptive runs exceeded the
allowance at all, and the summed propagator time was $6.2\%$ (fixed) and $7.4\%$
(adaptive) of the summed process time of all 40 runs. The share was much larger
in some short solved runs, up to about two fifths, and the untimed admission
tests are not included.

The adaptive variant obtained more accepted bound changes with fewer LPs: 308
against 112. At the root, both variants obtained 62 accepted changes, the
adaptive one with 336 LPs instead of 1150; most of its other changes came from
callbacks at tree nodes. Of the 470 bounds the adaptive variant proposed, SCIP applied 308;
of the fixed variant's 200 proposals, it applied 112. SCIP's tightening
interface declines a change that it does not regard as a sufficient
improvement; the archive does not record individual reasons. Witness screening
removed only two directions, one in each seed of
\texttt{kall\_circlespolygons\_c1p13}. The archive does not record the outcome of
individual feasibility checks; a check of a floating-point LP solution in exact
arithmetic fails whenever the solution violates an active row by any positive
amount, which limits how often such points can screen. Nonroot, number-zero callbacks consumed
360 of the fixed variant's LPs without producing an accepted change. Directional
LPs that produced no usable bound numbered 4 (fixed) and 26 (adaptive); their
individual statuses were not retained. All 48 unsupported callbacks of each arm
were on \texttt{qp3}, whose local box remained unbounded; the initially unbounded
\texttt{powerflow0009r} became eligible after SCIP tightened its bounds. The
numerical incumbent check never rejected SCIP's incumbent, and 124 (fixed) and
120 (adaptive) callbacks ran without a cutoff. The time limit was a stopping
criterion rather than an exact deadline: all 78 time-limited runs exceeded ten
seconds slightly, the longest solver call took 10.093 seconds, and the longest
process 10.606 seconds.

\subsection{What the comparison shows}\label{sec:experiments-interpretation}

The comparison supports four conclusions, each confined to this cohort,
budget, and implementation. Neither added policy changed which runs were solved
within ten seconds. The adaptive variant solved fewer auxiliary LPs than the
fixed variant but spent more time in the propagator. Neither added policy showed
a net speedup: on the runs that every arm solved, both took longer, by amounts
comparable to their own recorded propagator time. The improvement in mean public
gap comes from one model through a better incumbent; on the other public models
both policies had worse mean gap scores than native.

Several factors prevent stronger conclusions in either direction. The cohort
has twenty models, twelve of them selected using historical results and eight
synthetic, and the two seeds are repetitions rather than independent families.
The ten-second budget left 22 of 24 public runs unsolved in every arm, so PAR-2
differences reflect only the 14 solved runs. Most of these are short synthetic
runs, in which a fixed overhead of about half a second is a large share of the
process time and the recorded times move in steps of about 0.05 seconds. The
machine was shared with unrelated jobs. Under
these conditions the small timing differences do not establish a slowdown,
and the absence of new solves does not show that longer runs would also see
none. The propagator was written in Python and builds its relaxation in exact
rational arithmetic at every callback; its per-callback cost is a property of
this implementation, and an integrated implementation could have a very
different cost structure. The policy was compared as a whole, not by a factorial
design over ranking, screening, and the pilot. Finally, the protected-box,
cutoff-threshold, and residual certificates of
Sections~\ref{sec:certificates}--\ref{sec:residual} were not used by the measured
policy, so the comparison says nothing about their computational value.

The negative result matters for three reasons. First, it shows concretely that
the number of auxiliary LPs is not a cost measure. Work outside the LP solves
took about a third of the propagator time of the fixed variant and more than half
of that of the adaptive variant, and a policy that saved LPs by stopping early
spent the savings, and more, on additional callbacks. Second, the recorded
propagator time was at most $1.00501$ seconds per run and amounted to
$6.2\%$ and $7.4\%$ of the cohort's total
process time. A better selection rule can directly save at most the time the
propagator itself spends. A worthwhile gain must come from the search---fewer
nodes, earlier incumbents, or new solves. The search did change in places, as
\texttt{bayes2\_50} and the small bilinear cycles show, but 308 accepted
reductions produced no additional solve and no demonstrated overall speedup
within ten seconds. The archive records associations between policies and
search outcomes; without an ablation it does not isolate the effect of
particular reductions. This agrees with the precursor study, in which tightened
boxes reduced nodes but the final-solve savings did not offset the cost of
obtaining them. Third, it marks the boundary between the two kinds of results
in this paper. The certificates bound how much a specified operator can still
change a box or a relaxation bound, which can justify skipping work. They do not
predict that the work, if done, would save time, and Section~\ref{sec:effort}
shows that no budget statement can supply that prediction. On this evidence,
native SCIP remains the appropriate default, and the added policy is an
experimental implementation.
